\documentclass[a4paper,12pt]{article}

\usepackage[a4paper,margin=16mm]{geometry}
\usepackage[german,main=french]{babel}
\usepackage{fontspec}
\usepackage{microtype}

\setmainfont{EBGaramond}
[
  Extension      = .otf,
  UprightFont    = *-Regular,
  ItalicFont     = *-Italic,
  BoldFont       = *-Bold,
  BoldItalicFont = *-BoldItalic,
  Numbers        = Lowercase
]

\usepackage{paracol}
\newcommand\chunks[2]{%
\begin{leftcolumn*}\begin{otherlanguage}{german}%
{#1}%
\end{otherlanguage}\end{leftcolumn*}%
\begin{rightcolumn}\begin{otherlanguage}{french}%
{#2}%
\end{otherlanguage}\end{rightcolumn}%
}

\newcommand\strophe[1]{\subsection*{\normalsize\normalfont\itshape #1}}

\begin{document}\thispagestyle{empty}
\setlength{\columnsep}{8mm}


\begin{paracol}{2}

  \chunks{\section*{Durch Zärtlichkeit und Schmeicheln}}
         {}

  \chunks{\noindent\raggedright\small Wolfgang Amadeus Mozart, \textit{Die Entführung aus dem Serail}}
         {\noindent\raggedright\small Wolfgang Amadeus Mozart, \textit{L’Enlèvement au sérail}}

  \chunks{\noindent\raggedright\small Libretto von Johann Gottlieb Stephanie (1782)}
         {\noindent\raggedright\small Livret de Johann Gottlieb Stephanie (1782)}

  \chunks{\strophe{Zweiter Aufzug, erster Auftritt, Nr.~8, Arie}}
         {\strophe{Deuxième acte, première scène, n\textsuperscript{o}~8, air}}

  \chunks{\noindent\small(Garten am Palast des Bassa Selim; an der Seite Osmins Wohnung. Osmin, Blonde.)}
         {}

  \chunks{\begin{verse}
          \textsc{Blonde}\\[\medskipamount]
          Durch Zärtlichkeit und Schmeicheln,\\
          Gefälligkeit und Scherzen,\\
          Erobert man die Herzen\\
          Der guten Mädchen leicht:
          \end{verse}}
         {}

  \chunks{\begin{verse}
          Doch mürrisches Befehlen\\
          Und Poltern, Zanken, Plagen\\
          Macht, dass in wenig Tagen\\
          So Lieb’ als Treu entweicht.
          \end{verse}}
         {}
\end{paracol}
\end{document}
